\documentclass[conference]{IEEEtran}
\usepackage{amsmath,amssymb}
\usepackage{booktabs}
\usepackage{array}
\usepackage{url}
\title{A Task- and State-Conditioned Atlas of Coding Principles in Neuronal Spike Trains}
\author{\IEEEauthorblockN{Research Proposal}}
\begin{document}
\maketitle

\begin{abstract}
Neuronal spike trains are the observable output of individual neurons, yet the principles that map stimuli to spikes remain unresolved. Rate, latency, precise timing, synchrony, correlations, mixed selectivity, population geometry, and latent state can all carry information under some conditions. The central difficulty is that statistical availability, decodability, cross-condition generalization, and causal use by downstream circuits are different claims. This paper proposes a Task- and State-Conditioned Causal Neural Code Atlas (TSCC-A), a staged framework that compares these coding dimensions with matched null models and separates descriptive from functional evidence. A point-process encoding model combines stimulus filters, spike history, coupling, and state covariates. Nested decoders quantify incremental information from rate, timing, correlations, population structure, and geometry under held-out stimuli, states, tasks, neurons, and time windows. A future causal module selectively perturbs candidate timing patterns, correlation structure, or population subspaces while controlling marginal firing, movement, arousal, and nonspecific disruption. The proposal predicts a conditional atlas rather than a universal winner: a dimension is called available, generalizing, downstream-relevant, or functionally embedded only when it passes the corresponding evidence gate. The design is explicitly a proposal; no neural recordings, analyses, perturbations, or observed results are reported.
\end{abstract}

\begin{IEEEkeywords}
neural coding, spike trains, population coding, information theory, neural manifolds, causal inference, state dependence
\end{IEEEkeywords}

\section{Introduction}
The spike train is a compact but expressive interface between a neuron and the rest of the nervous system. It discards much of the continuous biophysics of membrane voltage while retaining a sequence of discrete events. A useful scientific question is therefore not merely whether a stimulus changes a firing rate. It is which properties of the event sequence carry information, at what temporal and spatial scale, under which behavioral conditions, and whether a downstream circuit uses those properties. The phrase "neural code" can refer to any of these levels, and collapsing them has produced apparently conflicting conclusions.

The rate-code view is attractive because spike counts are robust, interpretable, and often predictive. Temporal-code accounts point to latency, inter-spike intervals, and precise patterns. Population-code accounts emphasize that a stimulus can be represented by a distributed pattern even when single-neuron tuning is ambiguous. Correlation-based accounts ask whether the joint distribution of neurons carries information not present in their marginals. More recent work describes task-relevant population trajectories and low-dimensional manifolds. None of these descriptions is automatically the mechanism used by a biological circuit. A decoder can exploit a statistical regularity that a downstream synapse never reads.

This paper develops the selected Idea Agent direction, the Task- and State-Conditioned Causal Neural Code Atlas (TSCC-A). Its central hypothesis is:

\emph{If a coding dimension is functionally embedded in a neuronal population, its stimulus information should generalize across held-out stimuli and state conditions, improve a downstream-relevant decoder beyond matched nulls, and show a selective causal effect when the corresponding spike pattern or population subspace is perturbed.}

The hypothesis is strong enough to fail and broad enough to compare coding dimensions. It does not predict that one code dominates in every brain region, species, or task. Instead, it treats the code as a conditional object indexed by state, task, area, species, and analysis scale. This is important because the mapping between a stimulus and one neuron's response can change with attention, arousal, locomotion, internal state, task demand, adaptation, and network history.

The contributions are fourfold. First, TSCC-A defines a common representation ledger for rate, timing, correlations, population geometry, and hybrid features. Second, it combines encoding and decoding while enforcing held-out tests across stimulus and context. Third, it gives correlations and manifolds matched null models rather than treating them as intrinsically informative. Fourth, it establishes a promotion rule from information availability to functional embedding that requires causal validation. The resulting report is a research plan, not a claim that the atlas has been run.

\section{Background, Survey, and Research Gap}
\subsection{Spike trains as point-process observations}
For neuron $i$, let the observed spikes occur at event times $t_{ik}$. A distributional representation of the train is
\begin{equation}
s_i(t)=\sum_k\delta(t-t_{ik}).
\end{equation}
The delta representation is useful because it preserves event timing, while a count representation integrates it over a window. A rate statistic is consequently a transformation of the spike train, not an alternative physical output. The choice of bin width and window determines which temporal structure remains visible.

A conditional intensity provides a flexible statistical bridge from stimuli to events:
\begin{equation}
\lambda_i(t)=\exp\left[b_i+(k_i*x)(t)+\sum_j(w_{ij}*s_j)(t)+q_i(t)\right].
\end{equation}
Here $x(t)$ is the stimulus, $k_i$ is a stimulus filter, $w_{ij}$ describes spike history and coupling, and $q_i(t)$ contains measured state, task, movement, or arousal covariates. The model does not assert that the brain implements a log-linear GLM. It supplies a nested measurement language in which stimulus drive, temporal history, population coupling, and context can be compared.

\subsection{Rate and temporal coding}
The classic literature established that spike counts can transmit information about sensory variables and that the relevant temporal window is itself an empirical question [1], [2]. Rate is therefore a necessary baseline for a fair comparison. A rate-only model may predict a stimulus well when the nervous system integrates over tens or hundreds of milliseconds. It can also be useful for motor output, where a population firing rate correlates with a control variable.

However, a count discards event order. Temporal encoding studies formalized the possibility that latency, inter-spike intervals, and stimulus-dependent patterns carry information beyond the count [3]. Timing claims are especially sensitive to jitter, spike sorting, trial alignment, adaptation, and the chosen time scale. A short bin can make a rate change appear temporal; a long bin can erase a real timing code. The design must therefore compare count-matched and timing-preserving controls instead of relying on a fixed definition of temporal precision.

\subsection{Correlations and synchrony}
For a population, the joint response distribution can contain structure that is absent from each marginal response. Neural correlations can be redundant, synergistic, or detrimental to a particular decoder, depending on how signal correlations and noise correlations align [4]. Weak pairwise correlations can nevertheless produce strong collective states in a population model [5]. This result motivates explicit population likelihoods and careful scaling with the number of recorded units.

Complete-population visual recordings provided a concrete demonstration that spatio-temporal correlations can matter for neural signaling. A model that accounted for correlations extracted roughly 20 percent more visual information than an independent model in the reported setting and retained more information than an optimal linear decoder [6]. The result is important but conditional: it does not establish that correlation is always beneficial, nor that the same effect transfers to other populations, states, or tasks.

\subsection{Population codes and geometry}
In a population code, information can be distributed across many units with mixed selectivity. A neuron may respond to a conjunction of stimulus, task, and internal variables; this apparent complexity can make single-neuron labels unstable while preserving a population-level readout. Large-scale visual recordings revealed structured high-dimensional response geometry [7]. Work on movement has shown that neural manifolds and task-relevant trajectories can organize population activity in a compact dynamical description [8].

Geometry is a powerful summary, but it is not synonymous with mechanism. Dimensionality reduction can create a visually persuasive axis from noise, a latent variable can track a nuisance such as movement, and a stable subspace can be downstream-irrelevant. A geometry becomes stronger evidence when it predicts a relevant output under held-out conditions and when selective intervention on that geometry changes the output without globally disrupting activity.

\subsection{Information, encoding, and decoding}
Information-theoretic analysis provides tools for comparing response representations, but finite samples, estimator bias, binning, and population size matter [9]. An encoding model predicts spikes from a stimulus and context. A decoder predicts a stimulus or behavior from spikes. These are related but not symmetric tasks: multiple encoders can produce similar decoders, and a decoder can use correlations that are not biologically available to a downstream circuit.

The TSCC-A ledger records a condition-specific vector
\begin{equation}
\mathbf C(c)=\left(I_{\rm rate},I_{\rm time},I_{\rm corr},I_{\rm pop},I_{\rm state}\right),
\end{equation}
where $c$ includes state, task, area, and species. The components are incremental or absolute information estimates tied to matched baselines. The notation is not intended to rank dimensions on a universal axis; it records which dimension adds predictive value under a stated observation model.

\subsection{The unresolved mapping problem}
The stimulus-to-neuron relationship is not generally one-to-one. A response can be driven by stimulus features, recent spikes, neighboring neurons, movement, attention, internal state, or an unmeasured latent variable. Point-process analyses provide a way to model these effects jointly [10]. Yet many studies still report an isolated receptive field, an accuracy score, or a low-dimensional plot without a common cross-condition test. Three gaps follow.

First, rate, timing, correlation, and geometry are rarely compared with the same trials, capacity controls, and held-out conditions. Second, stimulus-neuron mappings are often evaluated within a state and task, leaving their transfer properties unclear. Third, information that can be decoded offline is frequently interpreted as information the brain uses, even though this requires a downstream or causal test. TSCC-A addresses these gaps as G1--G6: matched multidimensional benchmarking, conditional mappings, decodability versus causal use, sampling-dependent correlation scaling, geometry versus implementation, and temporal measurement controls.

\section{Research Questions and Planned Contributions}
The primary research question is: \emph{which coding dimensions of neuronal spike trains carry stimulus information under particular state and task conditions, and which of those dimensions are functionally used by downstream circuits?}

The study has five secondary questions:
\begin{enumerate}
\item How much predictive information is added by timing beyond count-matched rate?
\item When do correlations add information, and when do they mainly add redundancy or instability?
\item Do population-level geometries generalize when states, tasks, neurons, or stimuli are held out?
\item Which apparent stimulus-neuron mappings are context-bound?
\item What evidence is sufficient to call a dimension functionally embedded rather than merely decodable?
\end{enumerate}

The planned contributions are a preregistered representation ledger, a cross-condition generalization matrix, a correlation and timing null suite, a population-geometry relevance test, and a causal promotion rule. The atlas will retain null and unresolved outcomes rather than forcing a single label.

\begin{table}[t]
\caption{Evidence gates for a coding dimension}
\label{tab:gates}
\centering
\footnotesize
\begin{tabular}{p{0.19\columnwidth}p{0.70\columnwidth}}
\toprule
Status & Required evidence \\
\midrule
Available & Held-out stimulus information in at least one prespecified condition. \\
Generalizing & Information persists when state, task, stimulus, neuron, or time window is held out. \\
Downstream-relevant & Predicts a relevant behavior, decision, motor output, or cross-area response beyond nuisance controls. \\
Functionally embedded & Selective perturbation changes the relevant output beyond sham and matched nonspecific disruption. \\
\bottomrule
\end{tabular}
\end{table}

\section{Idea Source Checkpoints and Direction Selection Audit}
The Idea Agent began with the accepted Survey gaps and generated five routes. R1 decomposes rate and timing information. R2 evaluates correlation-aware population decoders. R3 analyzes state-conditioned neural manifolds. R4 creates a causal downstream assay. R5 combines the routes into TSCC-A. The route selection is not a claim that R5 is empirically superior; it is a coverage decision based on the scientific question.

R1 directly tests the rate-versus-timing dispute but leaves population geometry and use unresolved. R2 measures joint information and redundancy but remains a decoder result unless connected to a biological readout. R3 can explain why individual neuron mappings change while population behavior remains stable, but geometry alone is correlational. R4 has the strongest functional interpretation and the highest practical risk because selective perturbation, neural stability, ethics, and safety all matter. R5 makes R1--R4 nested stages, so partial completion remains informative while preserving a path to a stronger causal conclusion.

The Idea Agent's evolution checkpoints were:
\begin{enumerate}
\item The initial defect was an under-specified rate-versus-timing debate. The edit was to define matched representations and count-preserving nulls.
\item The second defect was confusing a high-performing population decoder with biological use. The edit was to add a downstream-relevance gate.
\item The third defect was treating a stimulus-neuron map as invariant. The edit was to condition models on state and task and require cross-condition tests.
\item The final defect was treating a manifold as causal. The edit was to add selective perturbation and nuisance-matched controls.
\end{enumerate}

The rejected portfolio member was a single universal code. The evidence map does not support that scope, and the proposition is too broad unless its context and time scale are specified. A decoder-only universal criterion was also rejected as a complete direction because it cannot distinguish availability from use. The primary direction is therefore TSCC-A, with rate-timing, correlation-population, and manifold analyses retained as competitive components.

\section{Problem Definition, Assumptions, and Hypotheses}
\subsection{Operational definitions}
Let $S$ denote the stimulus, $R$ the multineuron spike response, and $C$ the measured context. The context includes behavioral state, task, brain area, species, movement, arousal, and trial history where available. A stimulus-neuron mapping is the conditional model $p(R\mid S,C)$ rather than a fixed function from a stimulus label to a single cell. A coding dimension is a feature family that changes the predictive distribution or relevant decision when added to a nested model.

Rate is a count or smoothed intensity in a preregistered window. Timing is event order, latency, or inter-spike structure after count-matched controls. Correlation is joint dependence beyond marginal tuning, including noise correlations and synchrony. Population coding is distributed information across units. Geometry is the structure of trajectories, subspaces, and latent variables estimated without test-set leakage. A hybrid code combines dimensions and is not assumed to be reducible to their independent contributions.

\subsection{Assumptions}
The proposed analysis assumes that trials can be aligned to a stimulus and task epoch, that relevant nuisance variables can be measured or at least recorded as unknowns, and that the spike trains have sufficient quality for the temporal resolution being claimed. It assumes no particular species, area, recording modality, or decoder architecture. It assumes that a causal module, if later pursued, will occur only under appropriate approvals and will be interpreted with sham and nonspecific controls.

These assumptions are testable. If stimulus timing is uncertain, the timing component is marked unresolved. If population sampling is sparse, the atlas reports population-size uncertainty. If movement or arousal cannot be measured, claims about state dependence are limited. A limitation in measurement is not silently converted into evidence for a coding dimension.

\subsection{Hypotheses and falsification}
H1 predicts that rate will carry reproducible information in at least some conditions but will not dominate every held-out context. H2 predicts that temporal features will add information in conditions where precise event timing is stimulus-locked after count matching. H3 predicts that correlation effects will depend on signal-noise alignment and population size. H4 predicts that population geometry will provide a context-dependent, sometimes stable, representation of task variables. H5 predicts that individual stimulus-neuron mappings will change with context more often than a population-level relevant subspace. H6 predicts that only a subset of statistically decodable features will pass the downstream and causal gates.

The framework is falsified as a useful unified account if a rate-only model matches the full model across all held-out dimensions and all additional features fail matched-null, relevance, and causal tests. This would support a rate-dominant explanation for the sampled conditions, not a universal statement about all nervous systems. A timing gain that disappears under count-matched jitter, a correlation gain that disappears under correlation-preserving shuffles, or a manifold effect explained by movement is rejected as dimension-specific evidence.

\section{Study Design and Methods}
\subsection{D0: Preregistration and data contract}
Before analysis, freeze the stimulus set, state and task labels, neuron inclusion rule, temporal resolution, train-validation-test partitions, estimator family, nuisance covariates, null models, and multiple-comparison procedure. The data contract stores species, brain area, recording modality, sampling rate, spike-sorting confidence, trial exclusions, session identifiers, stimulus identity, behavioral state, task epoch, movement, arousal, and downstream output. The data schema also records missing fields so that a missing state label cannot be mistaken for a stable state.

The same trials and folds are used for all candidate representations. Hyperparameters and feature transforms are selected within training folds. Held-out splits include stimuli, states, tasks when available, neurons, and time windows. The primary outputs are effect sizes, proper scoring rules, calibration error, information estimates, uncertainty intervals, and scaling curves. Classification accuracy is secondary because it can hide calibration failure and class imbalance.

\subsection{D1: Stimulus-response encoding}
Fit nested point-process GLMs for each neuron and for the population. The rate-only model contains a stimulus filter and baseline. The timing model adds latency or temporal basis functions. The history model adds self-history. The coupling model adds neighboring spike histories. The full model adds state, task, movement, and arousal covariates. Regularization is chosen inside training folds, and hierarchical bootstrap accounts for trials, neurons, and sessions.

The outputs are receptive-field or encoding-filter estimates, latency distributions, firing-rate tuning, temporal precision, state-dependent modulation, residual coupling, and posterior or bootstrap uncertainty. The analysis asks whether the stimulus term remains predictive after context and history are included. It does not treat an improved likelihood as proof that the biological system uses that exact parametric form.

\subsection{D2: Multidimensional decoding}
Construct six decoders: count-based rate, temporal point-process, correlation-aware, latent-variable population, manifold, and hybrid. The rate decoder uses count features; the temporal decoder retains event times; the correlation decoder models joint dependencies; the population decoder uses distributed activity; the manifold decoder projects into a training-fold latent coordinate system; the hybrid decoder combines prespecified families. Capacity-matched baselines prevent a flexible model from winning solely through parameter count.

For an information comparison, the incremental temporal contribution is
\begin{equation}
\Delta I_{\mathrm{time}}=I(S;R_{\mathrm{full}})-I(S;R_{\mathrm{rate}}).
\end{equation}
The same construction estimates incremental correlation, population, and state contributions. Estimates are accompanied by finite-sample diagnostics, bootstrap intervals, and a record of binning or kernel choices. A decoder that predicts a stimulus from a train is labeled "available" under Table~\ref{tab:gates}; it is not labeled "used" without the later gates.

\subsection{D3: Null and robustness suite}
Count-matched temporal jitter preserves counts and coarse rate while disrupting fine event order. Trial shuffles preserve marginal tuning while disrupting trial-level coordination. Conditional-independence nulls remove joint structure after conditioning on stimulus and state. Correlation-preserving stimulus shuffles test whether a correlation effect depends on stimulus alignment rather than correlation magnitude alone. Population-size bootstrap curves reveal whether information gains persist as more neurons are sampled. State-label permutations test whether context effects exceed session drift. Spike-sorting and missing-unit sensitivity analyses test whether timing and correlation findings survive measurement uncertainty.

The nulls are not interchangeable. A jitter null addresses fine timing; an independent null addresses coupling; a population bootstrap addresses sampling. The analysis reports the comparison that failed as well as the one that passed. Multiple coding dimensions and conditions require a preregistered false-discovery procedure or hierarchical shrinkage. If the data cannot support a reliable estimator, the atlas returns an unresolved field instead of a directional claim.

\subsection{D4: State and task generalization}
Estimate the conditional mapping $p(R\mid S,C)$ and construct a generalization matrix. Train and test within state, across state, within task, across task, within area, and across area where the stimulus and output spaces overlap. Candidate context variables include arousal, locomotion, attention, task demand, internal state, and trial history. Hierarchical effects distinguish a common coding dimension from a context-specific gain.

The important prediction is conditional stability, not universal invariance. Individual neurons may remap while a population subspace remains predictive. Conversely, a stable receptive field in one session may fail across task demands. The matrix makes these alternatives visible and prevents a model trained in one condition from being presented as a general law.

\subsection{D5: Population geometry and downstream relevance}
Apply dimensionality reduction inside training folds and project held-out trials into the learned coordinates. Report participation ratio, effective dimensionality, trajectory separation, subspace alignment, and cross-condition stability. Predict behavior, decision, motor output, or a cross-area response with cross-validated models. Control for firing-rate gain, movement, trial history, arousal, session drift, and the number of units.

The downstream-relevance gate requires that a candidate coding feature predict the relevant output beyond a nuisance-only model and retain performance on held-out conditions. A low-dimensional trajectory that does not predict a relevant endpoint is descriptive. A high-dimensional distributed representation that predicts the endpoint is not rejected because it lacks a simple manifold. This protects the analysis from treating visual simplicity as biological necessity.

\subsection{D6: Future causal validation}
The causal module is deliberately separated from the current proposal artifact. If a future approved study can support it, choose one candidate dimension at a time. A timing intervention should alter event timing while preserving count as closely as possible. A correlation intervention should alter synchrony or joint structure while matching marginal rates. A geometry intervention should target a task-relevant population subspace while minimizing global activity changes. Each intervention requires sham stimulation, non-candidate population controls, nuisance-matched disruption, and a preregistered stopping rule.

For downstream output $\phi$, define
\begin{equation}
\tau_k=\mathbb E[\phi\mid do(u_k=1),c]-\mathbb E[\phi\mid do(u_k=0),c],
\end{equation}
where $u_k$ targets coding dimension $k$ and $c$ fixes stimulus, task, state, and nuisance controls. A selective effect supports functional embedding only if it exceeds matched nonspecific disruption, is replicated across held-out contexts, and cannot be explained by movement, arousal, recording quality, or overall excitability. A global performance collapse is not evidence for a specific code.

\begin{table*}[t]
\caption{TSCC-A analysis matrix}
\label{tab:matrix}
\centering
\footnotesize
\begin{tabular}{p{0.15\textwidth}p{0.23\textwidth}p{0.22\textwidth}p{0.28\textwidth}}
\toprule
Stage & Representation & Primary comparison & Interpretation if positive \\
\midrule
D1 & Rate and timing & Nested point-process models; count-matched temporal controls & Timing or rate adds predictive information under the sampled condition. \\
D2 & Correlation and synchrony & Joint model versus independent and correlation-preserving nulls & Joint structure contributes information beyond marginal tuning. \\
D2/D4 & Population and hybrid & Capacity-matched decoder across held-out stimuli and contexts & Distributed or combined features generalize under the stated context shift. \\
D5 & Manifold and subspace & Geometry versus nuisance-only behavior prediction & A population geometry is relevant to the measured output. \\
D6 & Candidate functional dimension & Selective perturbation versus sham and matched disruption & The dimension is causally implicated in the downstream endpoint. \\
\bottomrule
\end{tabular}
\end{table*}

\subsection{Failure handling and reproducibility}
The analysis stores raw-data identifiers, preprocessing versions, fold assignments, model specifications, seeds, estimator settings, null definitions, exclusion logs, and negative results. A failed schema or quality batch is discarded and marked \texttt{needs\_human\_input}; it is not filled with a model-generated positive. The final proposal retains an empty observed-results field until an approved future execution supplies measurements.

\section{Expected Outcome Branches and Conditional Conclusions}
The design has several scientifically meaningful outcomes. If rate dominates within and across held-out conditions, the conclusion is a rate-dominant code for those conditions, with timing and correlation contributions bounded by the uncertainty of the controls. This would be a strong, direct result, not a failure of the project. It would also identify where additional dimensions are unnecessary.

If timing adds information after count-matched jitter and the gain generalizes across state or task, the evidence supports a temporal component. If the gain is confined to one alignment or disappears under sorting sensitivity, it remains a measurement-dependent association. If correlations improve a population decoder and the effect survives correlation-preserving controls and population-size analysis, the result supports stimulus-aligned joint coding. If they only improve in-sample likelihood or disappear with a matched null, the effect is not promoted.

If individual neurons remap but a population trajectory generalizes, the atlas supports a distributed, context-conditioned representation. If the trajectory is stable only within a session or is explained by movement, the result describes a local geometry rather than an invariant code. If a high-dimensional representation predicts the output while manifold reduction does not, the atlas records a distributed downstream-relevant code without forcing low dimensionality.

The strongest branch occurs when one dimension passes all four gates: availability, generalization, downstream relevance, and selective causal effect. Even then, the correct conclusion is conditional: the dimension is functionally embedded for the tested species, area, task, state, and timescale. It is not evidence for a universal neural code. The weakest branch is a complete null in which rate-only coding matches the full model and all added dimensions fail. That result would narrow the code for the sampled conditions and test the framework's parsimony.

The study cannot answer whether a coding principle is universal without broad replication. It can answer a more actionable question: what feature of a spike train is predictive, transferable, relevant, and causally selectable under a defined biological context? This distinction is the main scientific value of the proposal.

\section{Risks, Limitations, and Human Review Requirements}
\subsection{Measurement and sampling}
Spike sorting errors can create false timing precision and distort correlations. Limited populations can make a distributed code appear absent. The protocol responds with sorting sensitivity, population-size curves, uncertainty intervals, and explicit quality fields. These controls reduce but cannot eliminate sampling limitations. Human review is required before promoting a result whose temporal scale approaches the recording or sorting resolution.

\subsection{Model and estimator dependence}
No decoder family is neutral. Flexible models can overfit, while simple models can miss nonlinear interactions. Cross-validation, capacity matching, proper scoring rules, and multiple representations reduce this risk. Information estimators are sensitive to sample size and discretization; statistical reviewers must inspect estimator diagnostics and the analysis code before publication.

\subsection{Context and causal interpretation}
State is partly latent. Measured arousal or movement may not capture all internal variables. A cross-state failure could reflect a hidden confound rather than a changed code. A perturbation can change many variables at once, including arousal, motor output, and recording quality. For this reason, the causal module is not a default execution step. It requires animal or human ethics review where applicable, a safety plan, data governance, humane endpoints for animal work, and independent interpretation of the perturbation controls.

\subsection{Scope and transfer}
The proposed atlas does not guarantee transfer across species, brain areas, or modalities. A result from a sensory population cannot be presented as a motor principle without a new test. Cross-species comparison requires harmonized task definitions and hierarchical uncertainty. The Survey retrieval audit was partial because the local OpenAlex and AnySearch connections were refused; established sources and a publisher-page verification were used, but a future systematic review should refresh the bibliography and screen new papers.

\subsection{Human review checklist}
Before any real execution, experts must review the data-use permissions, animal or human protocol, privacy and governance requirements, perturbation safety, power analysis, spike-sorting quality, estimator choice, stopping rules, and the distinction between observed and expected outcomes. The Author Agent preserves these requirements as explicit review fields rather than treating them as completed evidence.

\section{Conclusion}
Neuronal spike trains do not point to one context-free coding principle. They support a family of conditional hypotheses about rates, timing, coordination, population structure, geometry, and state. The proposed TSCC-A turns this pluralism into a testable workflow. It compares representations on common trials, uses nulls matched to the claim, tests generalization across state and task, and reserves the strongest functional label for selective causal evidence.

The central deliverable is therefore not a universal answer of "rate" or "timing." It is an evidence atlas that says which dimension is available, when it transfers, whether it predicts a relevant output, and whether downstream circuitry is selectively sensitive to it. That formulation preserves the explanatory power of modern population neuroscience while preventing an offline decoder, a single receptive field, or an attractive manifold from being mistaken for a demonstrated biological code.

\appendices
\section{Variable and Operational Definition Appendix}
\begin{table}[h]
\caption{Core variables}
\label{tab:variables}
\centering
\footnotesize
\begin{tabular}{p{0.25\columnwidth}p{0.61\columnwidth}}
\toprule
Symbol & Definition \\
\midrule
$S$ & Stimulus identity or continuous stimulus feature. \\
$R$ & Single-neuron or population spike response. \\
$C$ & State, task, area, species, movement, arousal, and history context. \\
$\lambda_i(t)$ & Conditional intensity of neuron $i$. \\
$\phi$ & Downstream behavior, decision, motor output, or cross-area endpoint. \\
$\tau_k$ & Controlled effect of perturbing candidate dimension $k$. \\
\bottomrule
\end{tabular}
\end{table}

The terms "available," "generalizing," "downstream-relevant," and "functionally embedded" are ordinal evidence statuses, not universal scores. A coding dimension can remain available without passing a later gate. "Unknown" means that the design does not identify the quantity under the available data; it does not mean zero.

\section{Evidence Coverage and Review Appendix}
The Survey source registry is organized by the claim it is allowed to support. E1--E3 support the rate, timing, and event-process vocabulary; their boundary is that a representation may carry information without being the universal code. E4--E6 support the conditional role of correlation and population statistics; their boundary is that correlation effects depend on stimulus alignment, population structure, and the decoder. E7--E8 support population geometry and neural-manifold descriptions; their boundary is that a latent coordinate or stable subspace is not automatically a causal channel. E9 supports information-theoretic and decoding methodology; its boundary is that decodability is weaker evidence than downstream use. E10 supports conditional point-process models; its boundary is that the model is a comparison language rather than a literal neural mechanism.

The current evidence status is "COMPLETE WITH RETRIEVAL LIMITATION." OpenAlex and AnySearch were attempted for review-oriented queries but the local connection was refused. The Nature publisher page for Pillow et al. was checked in the application browser for title, authorship, journal, volume, pages, abstract, and DOI metadata. No downstream section treats the retrieval limitation as evidence for or against any coding hypothesis.

\begin{thebibliography}{10}
\bibitem{rieke}
F. Rieke, D. Warland, R. de Ruyter van Steveninck, and W. Bialek, \emph{Spikes: Exploring the Neural Code}. Cambridge, MA, USA: MIT Press, 1997.

\bibitem{bialek}
W. Bialek, F. Rieke, R. de Ruyter van Steveninck, and D. Warland, \textquotedblleft Reading a neural code,\textquotedblright{} \emph{Science}, vol. 252, no. 5014, pp. 1854--1857, 1991.

\bibitem{theunissen}
F. Theunissen and J. P. Miller, \textquotedblleft Temporal encoding in nervous systems: A rigorous definition,\textquotedblright{} \emph{J. Comput. Neurosci.}, vol. 2, pp. 149--162, 1995.

\bibitem{averbeck}
J. H. Averbeck, P. E. Latham, and A. Pouget, \textquotedblleft Neural correlations, population coding and computation,\textquotedblright{} \emph{Nat. Rev. Neurosci.}, vol. 7, pp. 358--366, 2006.

\bibitem{schneidman}
E. Schneidman, M. J. Berry II, R. Segev, and W. Bialek, \textquotedblleft Weak pairwise correlations imply strongly correlated network states in a neural population,\textquotedblright{} \emph{Nature}, vol. 440, pp. 1007--1012, 2006.

\bibitem{pillow}
J. W. Pillow \emph{et al.}, \textquotedblleft Spatio-temporal correlations and visual signalling in a complete neuronal population,\textquotedblright{} \emph{Nature}, vol. 454, pp. 995--999, 2008.

\bibitem{stringer}
C. Stringer \emph{et al.}, \textquotedblleft High-dimensional geometry of population responses in visual cortex,\textquotedblright{} \emph{Nature}, vol. 571, pp. 361--365, 2019.

\bibitem{gallego}
C. Gallego \emph{et al.}, \textquotedblleft Neural manifolds for the control of movement,\textquotedblright{} \emph{Neuron}, vol. 94, pp. 978--984, 2017.

\bibitem{quiroga}
R. Q. Quiroga and S. Panzeri, \textquotedblleft Extracting information from neuronal populations: Information theory and decoding,\textquotedblright{} \emph{Nat. Rev. Neurosci.}, vol. 10, pp. 173--185, 2009.

\bibitem{truccolo}
G. Truccolo, U. T. Eden, M. R. Fellows, J. P. Donoghue, and E. N. Brown, \textquotedblleft A point process framework for relating neural spiking activity to spiking history, neural ensemble, and extrinsic covariate effects,\textquotedblright{} \emph{J. Neurophysiol.}, vol. 93, pp. 1074--1089, 2005.
\end{thebibliography}

\end{document}
